% ===================== فصل ۳ =====================
\chapter{عملگرها و عبارت‌ها}
\label{ch:operators}

تا اینجا یاد گرفتیم چگونه داده را در متغیرها نگه داریم. اما برنامه‌نویسی فقط
نگه‌داشتن داده نیست؛ بیشتر کار، \emph{محاسبه} و \emph{مقایسه} و
\emph{تصمیم‌گیری} روی داده‌هاست. ابزار این کارها \textbf{عملگر}ها هستند. در
این فصل با خانوادهٔ عملگرهای سلام آشنا می‌شویم.

\section{عبارت چیست؟}

هر چیزی که محاسبه شود و به یک \emph{مقدار} برسد، یک \textbf{عبارت}
(expression) است. \code{2 + 3} یک عبارت است که به \code{5} می‌رسد.
\code{سن > 18} هم یک عبارت است که به \code{درست} یا \code{نادرست} می‌رسد.
عملگرها همان نشانه‌هایی‌اند که این محاسبه‌ها را انجام می‌دهند: \code{+}،
\code{-}، \code{>} و مانند آن‌ها.

\section{عملگرهای حسابی}

برای کار با اعداد، این عملگرها را داریم:

\begin{center}
\begin{tabular}{c l l}
\toprule
\textbf{عملگر} & \textbf{کار} & \textbf{نمونه} \\
\midrule
\code{+} & جمع & \code{5 + 3}\;$\to$\;\code{8} \\
\code{-} & تفریق & \code{5 - 3}\;$\to$\;\code{2} \\
\code{*} & ضرب & \code{5 * 3}\;$\to$\;\code{15} \\
\code{/} & تقسیم & \code{10 / 3}\;$\to$\;\code{3} \\
\code{\%} & باقی‌ماندهٔ تقسیم & \code{10 \% 3}\;$\to$\;\code{1} \\
\code{**} & توان & \code{2 ** 10}\;$\to$\;\code{1024} \\
\bottomrule
\end{tabular}
\end{center}

دو عملگر آخر کمی نیاز به توضیح دارند.

\subsection{تقسیم صحیح در برابر تقسیم اعشاری}

به این برنامه برگرفته از مثال‌های واقعی سلام نگاه کنید:

\begin{salam}
تابع اصلی:
    الف اعشاری۶۴ = 7.0 / 2.0
    چاپ الف
    ب i32 = 7 / 2
    چاپ ب
پایان
\end{salam}

خروجی چنین است:

\begin{salamen}[language={}]
3.5
3
\end{salamen}

چرا \code{7 / 2} برابر \code{3} شد، نه \code{3.5}؟ زیرا وقتی \emph{هر دو} عدد
درست (صحیح) باشند، سلام تقسیم \emph{صحیح} انجام می‌دهد و بخش اعشاری را دور
می‌اندازد. اما \code{7.0 / 2.0} چون اعشاری است، نتیجهٔ دقیق \code{3.5} را
می‌دهد.

\begin{warn}
این یکی از رایج‌ترین سرچشمه‌های اشتباه برای تازه‌کارهاست. اگر می‌خواهید نتیجهٔ
تقسیم اعشاری باشد، دست‌کم یکی از دو عدد را اعشاری بنویسید، یا با
\code{بعنوان} تبدیلش کنید:
\begin{salam}
تابع اصلی:
    ب i32 = 7
    دقیق اعشاری۶۴ = (ب بعنوان اعشاری۶۴) / 2.0
    چاپ دقیق
پایان
\end{salam}
که چاپ می‌کند \code{3.5}.
\end{warn}

\subsection{باقی‌مانده (\%): قهرمان پنهان}

عملگر \code{\%} باقی‌ماندهٔ یک تقسیم را می‌دهد. \code{10 \% 3} برابر \code{1}
است، چون $10 = 3\times 3 + 1$. این عملگر ساده کاربردهای فراوان دارد؛
مهم‌ترینش تشخیص \emph{بخش‌پذیری} است: عددی بر ۲ بخش‌پذیر (زوج) است اگر
باقی‌ماندهٔ تقسیمش بر ۲ صفر باشد:

\begin{salam}
تابع اصلی:
    عدد صحیح۳۲ = 12
    اگر عدد % 2 == 0:
        چاپ "زوج است"
    وگرنه:
        چاپ "فرد است"
    پایان
پایان
\end{salam}

\begin{tip}
بازی معروف «فیزباز» که در فصل بعد می‌سازیم، تماماً بر پایهٔ همین عملگر
\code{\%} بنا شده است. هر جا با «هر چند تا یک‌بار»، «زوج/فرد» یا «بخش‌پذیری»
سر و کار داشتید، احتمالاً \code{\%} ابزار شماست.
\end{tip}

\section{عملگرهای مقایسه‌ای}

این عملگرها دو مقدار را مقایسه می‌کنند و پاسخشان همیشه \emph{منطقی} است ---
یعنی \code{درست} یا \code{نادرست}:

\begin{center}
\begin{tabular}{c l l}
\toprule
\textbf{عملگر} & \textbf{معنا} & \textbf{نمونه} \\
\midrule
\code{==} & برابر است با & \code{5 == 5}\;$\to$\;\code{درست} \\
\code{!=} & برابر نیست با & \code{5 != 3}\;$\to$\;\code{درست} \\
\code{<} & کوچک‌تر از & \code{3 < 5}\;$\to$\;\code{درست} \\
\code{>} & بزرگ‌تر از & \code{9 > 2}\;$\to$\;\code{درست} \\
\code{<=} & کوچک‌تر یا مساوی & \code{5 <= 5}\;$\to$\;\code{درست} \\
\code{>=} & بزرگ‌تر یا مساوی & \code{4 >= 7}\;$\to$\;\code{نادرست} \\
\bottomrule
\end{tabular}
\end{center}

\begin{warn}
به تفاوت \code{=} و \code{==} بسیار دقت کنید! یک \code{=} یعنی «این مقدار را
\emph{بگذار} در آن متغیر» (انتساب)، اما دو \code{==} یعنی «آیا این دو
\emph{برابرند}؟» (مقایسه). جابه‌جا گرفتن این دو، اشتباهی رایج است.
\end{warn}

\section{عملگرهای منطقی}

گاهی می‌خواهیم چند شرط را با هم ترکیب کنیم. عملگرهای منطقی این کار را می‌کنند:

\begin{center}
\begin{tabular}{c l l}
\toprule
\textbf{عملگر} & \textbf{معنا} & \textbf{درست است وقتی...} \\
\midrule
\code{\&\&} & و (and) & هر دو طرف درست باشند \\
\code{||} & یا (or) & دست‌کم یک طرف درست باشد \\
\code{!} & نقیض (not) & طرف مقابل نادرست باشد \\
\bottomrule
\end{tabular}
\end{center}

این برنامهٔ کوچک برگرفته از آزمون‌های واقعی سلام هر سه را نشان می‌دهد:

\begin{salam}
تابع اصلی:
    الف منطقی = درست
    ب منطقی = نادرست
    چاپ الف && ب, الف || ب, !الف
    چاپ 3 < 5, 5 <= 5, 9 > 2
پایان
\end{salam}

خروجی:

\begin{salamen}[language={}]
false true false
true true true
\end{salamen}

\begin{itemize}
  \item \code{الف \&\& ب} نادرست است، چون \code{ب} نادرست است و «و» نیاز دارد
        که \emph{هر دو} درست باشند.
  \item \code{الف || ب} درست است، چون دست‌کم یکی (\code{الف}) درست است.
  \item \code{!الف} نادرست است، چون \code{الف} درست بود و نقیضش نادرست می‌شود.
\end{itemize}

\begin{tip}
یک کاربرد روزمره: «اگر هوا آفتابی است \emph{و} تعطیل است، به پارک برویم.» این
جمله در سلام چنین می‌شود:
\begin{salam}
اگر آفتابی && تعطیل:
    چاپ "به پارک برویم"
پایان
\end{salam}
\end{tip}

\section{پیوند رشته‌ها}

عملگر \code{+} تنها برای اعداد نیست؛ روی رشته‌ها هم کار می‌کند و دو متن را به
هم \emph{می‌چسباند}. به این مثال واقعی نگاه کنید:

\begin{salam}
تابع اصلی:
    نام رشته = "سلام"
    سلام رشته = نام + " دنیا"
    چاپ سلام
    چاپ نام.len()
    چاپ سلام.len()
پایان
\end{salam}

خط \code{نام + " دنیا"} دو رشته را به هم چسباند و حاصل «سلام دنیا» شد. همچنین
\code{نام.len()} طول رشته (تعداد نویسه‌ها) را می‌دهد. در فصل ۸ بسیار بیشتر
دربارهٔ کار با رشته‌ها خواهیم گفت.

\begin{note}
عملگر \code{+} وقتی هر دو طرف عدد باشند، جمع می‌کند؛ و وقتی رشته باشند،
می‌چسباند. سلام از روی \emph{نوع} طرفین تشخیص می‌دهد کدام کار را انجام دهد.
\end{note}

\section{ترتیب اولویت عملگرها}

وقتی چند عملگر در یک عبارت باشند، سلام به ترتیبی مشخص آن‌ها را محاسبه می‌کند ---
درست مانند ریاضیات مدرسه. برای نمونه در \code{2 + 3 * 4}، نخست ضرب انجام
می‌شود و سپس جمع، پس نتیجه \code{14} است نه \code{20}. ترتیب کلی (از زودتر به
دیرتر) چنین است:

\begin{enumerate}
  \item پرانتز: \code{(\ldots)}
  \item نقیض و منفی یکانی: \code{!}، \code{-}
  \item توان: \code{**}
  \item ضرب، تقسیم، باقی‌مانده: \code{*}، \code{/}، \code{\%}
  \item جمع و تفریق: \code{+}، \code{-}
  \item مقایسه‌ها: \code{<}، \code{>}، \code{<=}، \code{>=}
  \item برابری: \code{==}، \code{!=}
  \item و: \code{\&\&}
  \item یا: \code{||}
\end{enumerate}

\begin{tip}
لازم نیست این جدول را حفظ کنید! هر وقت شک داشتید، از \textbf{پرانتز} استفاده
کنید. \code{(2 + 3) * 4} هیچ ابهامی ندارد و خواناتر هم هست. برنامه‌نویسان
باتجربه نیز برای روشنی بیشتر، فراوان از پرانتز بهره می‌برند.
\end{tip}

به این مثال واقعی از تبدیل دما دقت کنید که در آن پرانتز نقش کلیدی دارد:

\begin{salam}
تابع اصلی:
    چاپ "سلسیوس به فارنهایت"
    برای متغیر س = 0; س <= 40; س = س + 10:
        ف اعشاری۶۴ = (س بعنوان اعشاری۶۴) * 9.0 / 5.0 + 32.0
        چاپ س, "س =", ف, "ف"
    پایان
پایان
\end{salam}

اینجا پرانتز \code{(س بعنوان اعشاری۶۴)} تضمین می‌کند که نخست تبدیل نوع انجام
شود و سپس ضرب و تقسیم. (فرمول تبدیل سلسیوس به فارنهایت
$F = C \times \frac{9}{5} + 32$ است.)

\section{عملگرها روی انواع خودمان}

نکتهٔ شگفت‌انگیز سلام این است که می‌توانید عملگرها را برای \emph{نوع‌های
ساختهٔ خودتان} هم تعریف کنید. مثلاً اگر یک «نقطه» در صفحه بسازید، می‌توانید
معنای جمع دو نقطه را خودتان مشخص کنید:

\begin{salam}
ساختار نقطه:
    x i32 = 0
    y i32 = 0
    تابع عملگر+(دیگری نقطه) نقطه:
        بازگشت نقطه { x = این.x + دیگری.x, y = این.y + دیگری.y }
    پایان
پایان
تابع اصلی:
    الف نقطه = نقطه { x = 3, y = 4 }
    ب نقطه = نقطه { x = 1, y = 2 }
    ج نقطه = الف + ب
    چاپ ج.x, ج.y
پایان
\end{salam}

اینجا با تعریف \code{تابع عملگر+}، گفتیم که جمع دو نقطه یعنی جمع مؤلفه‌های
\code{x} و \code{y}. حالا \code{الف + ب} معنا پیدا می‌کند و خروجی \code{4 6}
می‌شود. نگران نباشید اگر هنوز \code{ساختار} را نمی‌شناسید؛ در فصل ۹ به‌تفصیل
به آن می‌پردازیم. این مثال فقط برای آن است که قدرت سلام را زود ببینید.

\begin{deep}[نگاهی به درون: عملگرها چگونه گسترش می‌یابند]
وقتی \code{الف + ب} را روی دو «نقطه» می‌نویسید، سلام در زمان کامپایل آن را به
فراخوانی تابع \code{عملگر+} تبدیل می‌کند چیزی شبیه
\code{الف.عملگر+(ب)}. هیچ هزینهٔ اضافه‌ای در زمان اجرا ندارد؛ این یک ترجمهٔ
ساده در زمان ساخت برنامه است. به همین شیوه می‌توانید \code{==}، \code{*} و
بسیاری عملگر دیگر را نیز برای نوع‌های خود تعریف کنید (فصل ۹).
\end{deep}

\section{خلاصهٔ فصل}

\begin{itemize}
  \item عبارت، هر چیزی است که به یک مقدار می‌رسد.
  \item عملگرهای حسابی: \code{+ - * / \% **}.
  \item تقسیم دو عدد صحیح، صحیح است؛ برای نتیجهٔ اعشاری، یک طرف را اعشاری
        کنید.
  \item عملگرهای مقایسه‌ای پاسخ منطقی می‌دهند.
  \item عملگرهای منطقی (\code{\&\&}، \code{||}، \code{!}) شرط‌ها را ترکیب
        می‌کنند.
  \item \code{+} روی رشته‌ها آن‌ها را به هم می‌چسباند.
  \item هنگام شک، از پرانتز برای روشن‌کردن ترتیب استفاده کنید.
\end{itemize}

\section{تمرین‌ها}

\exercise برنامه‌ای بنویسید که محیط و مساحت یک مستطیل با طول ۸ و عرض ۵ را
حساب و چاپ کند.

\exercise با استفاده از عملگر \code{\%}، برنامه‌ای بنویسید که بگوید عدد ۲۷ بر
۳ بخش‌پذیر است یا نه.

\exercise حاصل \code{2 + 3 * 4} و \code{(2 + 3) * 4} را در دو خط چاپ کنید و
تفاوت را ببینید.

\exercise دو متغیر منطقی بسازید (مثلاً «گرسنه» و «پول‌دارم») و با \code{\&\&}
تصمیم بگیرید که آیا باید غذا بخرید یا نه.

\exercise نام کوچک و نام خانوادگی خود را در دو رشتهٔ جداگانه بگذارید، آن‌ها
را با یک فاصله به هم بچسبانید و نام کامل را چاپ کنید. سپس طول نام کامل را
نیز چاپ کنید.
